% !PRINTED 316
\weberpage{316}
The two substitutions $1, \gamma_0$ form a group $Q_0$. Any substitution $\gamma_1$ of $Q$ that does not occur in $Q_0$ carries 0, 1 into two elements different from them, say 2, 3, and by the two substitutions $Q_0\gamma_1$ the pair 0, 1 passes into 2, 3 or into 3, 2. If, further, by $\gamma_2$ the pair 0, 1 passes into a third pair 4, 5, then 4, 5 is different both from 0, 1 and from 2, 3, and so we can decompose $Q$ into the subsidiary groups:

\[
Q = Q_0 + Q_0\gamma_1 + Q_0\gamma_2 + Q_0\gamma_3,
\]

whereby the eight elements (1) are decomposed into four pairs

\weq{2}{0, 1;\quad 2, 3;\quad 4, 5;\quad 6, 7}

so that by the two substitutions of one of the subsidiary groups $Q_0\gamma_i$ the pair 0, 1 passes into one of these four pairs.

Now if $k$ is any substitution from the group

\[
P = T_{24}\,Q,
\]

then the elements $\tau_0, \tau_1, \tau_2, \tau_3$ from $T_{24}$ and $\gamma_0', \gamma_1', \gamma_2', \gamma_3'$ from $Q$ can be so determined that

\[
k = \tau_0\gamma_0', \quad \gamma_1 k = \tau_1\gamma_1', \quad \gamma_2 k = \tau_2\gamma_2', \quad \gamma_3 k = \tau_3\gamma_3'
\]

or

\[
k = \tau_0\gamma_0' = \gamma_1^{-1}\tau_1\gamma_1' = \gamma_2^{-1}\tau_2\gamma_2' = \gamma_3^{-1}\tau_3\gamma_3',
\]

and this representation shows that by $k$ any one of the pairs (2) is carried into a pair of the same system; for, e.\,g., by $\gamma_1^{-1}$ the pair 2, 3 is carried into 0, 1, by $\tau_1$ the pair 0, 1 remains unchanged, and by $\gamma_1'$, which belongs to $Q$, the pair 0, 1 is carried into one of the pairs (2).

The group $T_{24}\,Q$ is therefore imprimitive, and indeed in such a way that we have four systems of imprimitivity. The roots 0, 1 satisfy a quadratic equation whose coefficients depend on the roots of a biquadratic equation.

Let us now pass to the case 2), in which the elements of a triple 1, 2, 3 and of a quadruple 4, 5, 6, 7 are each transitively connected among themselves by $T_{24}$. The eight quantities (1) here split into two quadruples:

\weq{3}{0, 1, 2, 3;\quad 4, 5, 6, 7,}

and by the substitutions of $Q$ the quantities of one of these quadruples are either only interchanged among themselves, or they are carried into the quantities of the other quadruple.
